\documentclass{article}
\usepackage{amsmath}
\title{When a live packet history forces framing overhead}
\author{}\date{}
\begin{document}\maketitle

\paragraph{Law and kernels.}
A rescue receiver should accept two pieces only when they discharge the
same authorized generation, not merely because the combined bytes parse.
This matters for per-inspection receipts even when two inspections produce
identical data. Each piece arrives from a predeclared positional source.
The trusted local inputs are origin, payload, and the declared cut. A
forgotten generation is not an input. Both histories in the comparison
actually reach the selected native SubT callbacks. Whether the second piece
was prepared by the same generation is an evaluator annotation; it is
not leaked to the controller. No forged packet, attacker, hardware
network, or deployment is assumed. The issue is delayed ordinary work.

\paragraph{Same observations, different settlement.}
In the first history the two pieces are from the new generation. In the
second the right arrives from a still-live older generation with the same
source, reception time and byte content. Under the same receiver observation
cut, doing anything other than accept forfeits valid on-time rescue in the
first history; acceptance would falsely settle the new generation in the
second. Every deterministic receiver has the same observation and hence the
same decision. This is not an obstruction for a receiver entitled to
identify equivalent content independent of generation, or one that already
has an enforced single-generation channel. A fresh scheduling record does
not provide that enforcement unless it excludes old prepared work reaching
the same acceptance cut.

\paragraph{Exact framing condition for declared equality receivers.}
Take finitely many eligible generations. Join two generations by an
undirected edge when the admitted source permits a cross-pair of original
fragments from them in a receiver cohort where both same-tag forms would
be accepted as one. The coloring law is restricted to codes with one
static tag per generation, a tag on both halves, and an equality test in
that cohort. A valid such assignment is exactly a proper coloring:
an edge with equal tags admits its witnessed cross-pair; different tags
reject every such pair. A same-generation pair uses the same tag by
construction. This gives a lower bound from simultaneously indistinct
candidates when they conflict pairwise, not a bound against every code,
time slot, or protocol. Cohort-wise exclusivity can be obtained without
payload bytes if enforced time windows, scoped authority, or external
qualified recognition already separate all mixed-generation histories.
Peak cohort size alone is insufficient for arbitrary static assignments:
three pair cohorts on three generations give a triangle requiring three
tags despite each allowing only two in its local cohort.

\paragraph{Coupled resource consequence.}
The headless replay sends ordinary one- and two-byte fragments at the same
cut through the pinned broker and captures the actual terminal callback
pairs. With no preexisting charges, a one-byte tag on each fragment
distinguishes the two histories and accepts only the coherent one. In the
matched remaining-credit case each sender has two of its original five
bytes left at the exact expiry boundary. The zero-byte framing handler
receives both and is ambiguous. With an added byte on each fragment, the
two-byte left fragment cannot arrive in time: the total available four
bytes are fewer than the five required, so no repartition of the
three-byte payload fixes this. The earlier zero-overhead rescue result
therefore depended on its single-generation law. A known exhaustively
valid expiry law, existing metadata channel, or larger reserve can change
the conclusion. Sending duplicate fragments without the missing
distinction does not. Native origin labels, staging provenance and payload
header are separate; the tested receiver is authored, not a claim about
a stock reassembly feature or about NeBula deployment.
\end{document}
